\documentclass[11pt,a4paper]{article}

% Packages
\usepackage[utf8]{inputenc}
\usepackage[T1]{fontenc}
\usepackage{lmodern}
\usepackage[margin=1in]{geometry}
\usepackage{graphicx}
\usepackage{hyperref}
\usepackage{xcolor}
\usepackage{listings}
\usepackage{booktabs}
\usepackage{tabularx}
\usepackage{longtable}
\usepackage{amsmath}
\usepackage{amssymb}
\usepackage{bytefield}
\usepackage{fancyhdr}
\usepackage{titlesec}
\usepackage{enumitem}
\usepackage{tcolorbox}
\usepackage{float}
\usepackage{caption}
\usepackage{subcaption}
\usepackage{multicol}

% Color definitions
\definecolor{skyblue}{RGB}{0,102,204}
\definecolor{darkblue}{RGB}{0,51,102}
\definecolor{codebg}{RGB}{245,245,245}
\definecolor{codegreen}{RGB}{0,128,0}
\definecolor{codegray}{RGB}{128,128,128}
\definecolor{codepurple}{RGB}{128,0,128}
\definecolor{warnorange}{RGB}{255,140,0}
\definecolor{alertred}{RGB}{204,0,0}

% Hyperref setup
\hypersetup{
    colorlinks=true,
    linkcolor=darkblue,
    urlcolor=skyblue,
    citecolor=darkblue,
    pdftitle={UAVLink Protocol v1.2 -- Official Documentation},
    pdfauthor={Krishnashis Das},
}

% Code listing style
\lstdefinestyle{cstyle}{
    backgroundcolor=\color{codebg},
    basicstyle=\ttfamily\footnotesize,
    breaklines=true,
    captionpos=b,
    commentstyle=\color{codegreen},
    keywordstyle=\color{skyblue}\bfseries,
    stringstyle=\color{codepurple},
    numberstyle=\tiny\color{codegray},
    numbers=left,
    numbersep=5pt,
    frame=single,
    rulecolor=\color{codegray},
    language=C,
    showstringspaces=false,
    tabsize=4,
    morekeywords={uint8_t,uint16_t,uint32_t,int8_t,int16_t,int32_t,bool,size_t,ssize_t,true,false},
}
\lstset{style=cstyle}

% Header/Footer
\setlength{\headheight}{14pt}
\pagestyle{fancy}
\fancyhf{}
\fancyhead[L]{\textcolor{darkblue}{\textbf{UAVLink Protocol v1.2}}}
\fancyfoot[C]{\thepage}
\renewcommand{\headrulewidth}{0.4pt}
\renewcommand{\footrulewidth}{0.4pt}

% tcolorbox environments
\newtcolorbox{notebox}{colback=blue!5!white,colframe=skyblue,title=\textbf{Note},fonttitle=\bfseries}
\newtcolorbox{warnbox}{colback=orange!5!white,colframe=warnorange,title=\textbf{Warning},fonttitle=\bfseries}
\newtcolorbox{alertbox}{colback=red!5!white,colframe=alertred,title=\textbf{Security},fonttitle=\bfseries}

% Title formatting
\titleformat{\section}{\Large\bfseries\color{darkblue}}{\thesection}{1em}{}
\titleformat{\subsection}{\large\bfseries\color{skyblue}}{\thesubsection}{1em}{}
\titleformat{\subsubsection}{\normalsize\bfseries\color{darkblue}}{\thesubsubsection}{1em}{}

\begin{document}

% ============================================================
% TITLE PAGE
% ============================================================
\begin{titlepage}
    \centering
    \vspace*{2cm}
    {\Huge\bfseries\textcolor{darkblue}{UAVLink Protocol}\par}
    \vspace{0.5cm}
    {\LARGE\textcolor{skyblue}{Version 1.2}\par}
    \vspace{0.3cm}
    {\large\textcolor{codegray}{Technical Documentation}\par}
    \vspace{2cm}
    {\large A High-Performance, Encrypted Binary Communication Protocol\\for Unmanned Aerial Vehicle Systems\par}
    \vspace{3cm}
    {\normalsize
    \textbf{Theorized and Developed by:} Krishnashis Das\\
    GitHub: \href{https://github.com/I-am-Krish}{\texttt{I-am-Krish}}\\[1em]
    \textbf{Co-Developed by:} Aadarsh Sinha\\
    GitHub: \href{https://github.com/Monarch666}{\texttt{Monarch666}}\\[2em]
    \textbf{Release Date:} March 2026\\
    \textbf{Document Revision:} 1.2.0\\
    \textbf{License:} Open Source
    }
    \vfill
    {\footnotesize\textcolor{codegray}{This document is the authoritative specification for the UAVLink protocol.\\All implementations must conform to the structures and algorithms described herein.}}
\end{titlepage}

\tableofcontents
\newpage

% ============================================================
% 1. INTRODUCTION
% ============================================================
\section{Introduction}

\subsection{What is UAVLink?}

\textbf{UAVLink} is a purpose-built, high-performance binary communication protocol designed for reliable, secure, and efficient data exchange between Unmanned Aerial Vehicles (UAVs) and Ground Control Stations (GCS). It operates over any byte-oriented transport layer---including UART serial links, UDP/IP networks, and radio telemetry modems---providing a unified framing, encryption, and integrity-checking layer for all UAV telemetry, command, and control traffic.

UAVLink addresses the critical shortcomings of existing UAV communication protocols:
\begin{itemize}[noitemsep]
    \item \textbf{Minimal overhead:} Bit-packed headers as small as 8 bytes minimize bandwidth consumption on constrained radio links.
    \item \textbf{Built-in security:} ChaCha20-Poly1305 AEAD encryption with 128-bit MAC tags protects against eavesdropping and tampering.
    \item \textbf{Real-time performance:} Deterministic O(1) memory allocation and zero-copy parsing for sub-millisecond latency.
    \item \textbf{Bandwidth efficiency:} Delta encoding, selective encryption, and message batching achieve up to 82.8\% bandwidth reduction.
    \item \textbf{Reliability:} CRC-16 integrity checks, replay protection, and forward error correction handle lossy links.
\end{itemize}

\subsection{Why ``UAVLink''?}

The name \textbf{UAVLink} was chosen to evoke the protocol's core mission: establishing a reliable, secure communication \emph{link} between systems operating in the \emph{sky} and those on the ground. It is short, memorable, and domain-specific to aerial robotics.

\subsection{Design Philosophy}

UAVLink was designed with the following principles:
\begin{enumerate}[noitemsep]
    \item \textbf{Efficiency First:} Every bit in the header is purposefully allocated. No wasted bytes.
    \item \textbf{Security by Default:} Encryption is not bolted on---it is integral to the packet format.
    \item \textbf{Determinism:} All operations have bounded, predictable execution time for hard real-time systems.
    \item \textbf{Layered Optimization:} Three phases of optimization can be adopted incrementally.
    \item \textbf{Portability:} Pure C99 with no external dependencies beyond the bundled Monocypher library.
\end{enumerate}

\subsection{Key Achievements}

\begin{table}[H]
\centering
\caption{UAVLink Performance Summary}
\begin{tabular}{@{}llll@{}}
\toprule
\textbf{Metric} & \textbf{Baseline} & \textbf{Optimized} & \textbf{Improvement} \\
\midrule
Bandwidth       & 3.68 kbps & 0.63 kbps & 82.8\% reduction \\
Parse Speed     & 250 $\mu$s & 125 $\mu$s & 2$\times$ faster \\
Crypto Speed    & 200 $\mu$s & 50 $\mu$s  & 4$\times$ faster (ARM NEON) \\
Memory Alloc    & 50 $\mu$s  & $<$1 $\mu$s & 50$\times$ faster (O(1) pool) \\
Total Latency   & 500 $\mu$s & 176 $\mu$s & 2.8$\times$ faster \\
\bottomrule
\end{tabular}
\end{table}

% ============================================================
% 2. ARCHITECTURE OVERVIEW
% ============================================================
\section{Architecture Overview}

\subsection{Protocol Stack}

UAVLink is organized into three optimization phases, each building upon the previous:

\begin{description}[style=nextline]
    \item[Core Protocol (Phase 0)] Packet framing, bit-packed headers, CRC-16 integrity, ChaCha20-Poly1305 AEAD encryption, byte-by-byte state machine parser, message serialization/deserialization, and nonce management.
    \item[Phase 1 --- Quick Wins] Selective encryption policies (60\% bandwidth reduction), crypto context caching (30\% speedup), and message batching (18\% bandwidth reduction).
    \item[Phase 2 --- Performance] Zero-copy parser (2$\times$ speed), fixed-size memory pool with O(1) allocation, and runtime hardware crypto detection (ARM NEON, x86 SSE2/AVX2).
    \item[Phase 3 --- Advanced] LZ4-style payload compression, Reed-Solomon forward error correction, and delta encoding for telemetry streams (57\% GPS bandwidth savings).
\end{description}

\subsection{File Organization}

\begin{table}[H]
\centering
\caption{Source File Map}
\begin{tabularx}{\textwidth}{@{}lX@{}}
\toprule
\textbf{File} & \textbf{Purpose} \\
\midrule
\texttt{uavlink.h / .c}           & Core protocol: headers, CRC, AEAD, parser, serialization, batching \\
\texttt{uavlink\_fast.h / .c}     & Phase 2: zero-copy parser, memory pool, hardware crypto detection \\
\texttt{uavlink\_compress.h / .c} & Phase 3: LZ4 compression, Reed-Solomon FEC, delta encoding \\
\texttt{uavlink\_hw\_crypto.h / .c} & Hardware-accelerated ChaCha20 (ARM NEON, x86 SSE2/AVX2) \\
\texttt{monocypher.h / .c}        & Portable ChaCha20-Poly1305 cryptography library \\
\texttt{uav\_simulator.c}         & Bidirectional UAV simulator (telemetry + command processing + ACKs) \\
\texttt{gcs\_receiver.c}          & Bidirectional GCS (telemetry display + interactive command menu) \\
\texttt{uavlink\_benchmark.c}     & Comprehensive performance profiler \\
\bottomrule
\end{tabularx}
\end{table}

% ============================================================
% 3. PACKET FORMAT
% ============================================================
\section{Packet Format Specification}

\subsection{Overall Packet Structure}

Every UAVLink packet follows this layout:

\begin{center}
\begin{tabular}{|c|c|c|c|c|}
\hline
\textbf{Base Header} & \textbf{Extended Header} & \textbf{Payload} & \textbf{MAC Tag*} & \textbf{CRC-16} \\
4 bytes & 4--13 bytes & 0--4095 bytes & 16 bytes* & 2 bytes \\
\hline
\end{tabular}
\end{center}

\begin{notebox}
The 16-byte Poly1305 MAC tag is present \emph{only} when the encrypted flag is set in byte 3. Unencrypted packets omit both the nonce (in the extended header) and the MAC tag, making them 24 bytes shorter.
\end{notebox}

\textbf{Packet size range:}
\begin{itemize}[noitemsep]
    \item Minimum: 10 bytes (empty payload, no encryption, no fragmentation, broadcast)
    \item Maximum: 4{,}122 bytes (4095-byte payload + full headers + MAC + CRC)
    \item Typical: 26--50 bytes (common telemetry messages)
\end{itemize}

\subsection{Base Header (4 Bytes)}

The base header is densely bit-packed to minimize overhead. Every bit serves a purpose.

\subsubsection{Byte 0: Start of Frame (SOF)}
Fixed value \texttt{0xA5}. The parser uses this as a synchronization marker to detect the beginning of a new frame in a continuous byte stream.

\subsubsection{Byte 1: Length[11:8] | Priority | StreamType[3:2]}
\begin{itemize}[noitemsep]
    \item \textbf{Bits 7--4:} Upper 4 bits of the 12-bit payload length field.
    \item \textbf{Bits 3--2:} 2-bit priority level (see Section~\ref{sec:priority}).
    \item \textbf{Bits 1--0:} Upper 2 bits of the 4-bit stream type field.
\end{itemize}

\subsubsection{Byte 2: StreamType[1:0] | Length[7:2]}
\begin{itemize}[noitemsep]
    \item \textbf{Bits 7--6:} Lower 2 bits of the stream type.
    \item \textbf{Bits 5--0:} Middle 6 bits of the payload length.
\end{itemize}

\subsubsection{Byte 3: Length[1:0] | Encrypted | Fragmented | Sequence[11:10]}
\begin{itemize}[noitemsep]
    \item \textbf{Bits 7--6:} Lower 2 bits of the payload length.
    \item \textbf{Bit 3:} Encrypted flag (1 = ChaCha20-Poly1305 AEAD applied).
    \item \textbf{Bit 2:} Fragmented flag (1 = packet is part of a multi-fragment message).
    \item \textbf{Bits 1--0:} Upper 2 bits of the 12-bit sequence number.
\end{itemize}

\begin{notebox}
The 12-bit payload length field supports payloads from 0 to 4095 bytes. The 12-bit sequence number rolls over at 4095 and is used for packet ordering and replay detection.
\end{notebox}

\subsection{Extended Header (4--13 Bytes)}

The extended header immediately follows the base header and contains routing, identification, and security fields.

\subsubsection{Always Present (4 Bytes)}

\textbf{Bytes 0--1: Sequence[9:0] | System ID}
\begin{itemize}[noitemsep]
    \item Upper 10 bits: Remaining sequence number bits (combined with base header for full 12-bit sequence).
    \item Lower 6 bits: Source system ID (0--63). Identifies the sending UAV or GCS.
\end{itemize}

\textbf{Bytes 2--3: Component ID | Message ID}
\begin{itemize}[noitemsep]
    \item Upper 4 bits: Component ID (0--15). Differentiates components on the same system (autopilot, gimbal, companion computer).
    \item Lower 12 bits: Message ID (0--4095). Defines the payload structure and semantics.
\end{itemize}

\subsubsection{Conditional: Target System ID (1 Byte)}
Present only for \texttt{CMD} and \texttt{CMD\_ACK} stream types. Contains the 6-bit target system ID (0 = broadcast to all systems).

\subsubsection{Conditional: Fragmentation Info (2 Bytes)}
Present only when the fragmented flag is set:
\begin{itemize}[noitemsep]
    \item \textbf{Fragment Index} (1 byte): Zero-based index of this fragment.
    \item \textbf{Fragment Total} (1 byte): Total number of fragments in the complete message.
\end{itemize}

\subsubsection{Conditional: Nonce (8 Bytes)}
Present only when the encrypted flag is set. Contains the 64-bit nonce used for ChaCha20-Poly1305 AEAD:
\begin{itemize}[noitemsep]
    \item \textbf{Bytes 0--3:} 32-bit monotonically increasing counter (little-endian). Ensures uniqueness.
    \item \textbf{Bytes 4--7:} 32-bit cryptographically random value. Adds entropy to prevent prediction.
\end{itemize}

\subsection{Payload (0--4095 Bytes)}

The payload contains the serialized message data. If encryption is enabled, this section contains \emph{ciphertext} (the plaintext encrypted with ChaCha20). The payload format depends on the message ID (see Section~\ref{sec:messages}).

\subsection{MAC Tag (16 Bytes, Conditional)}

Present only when the encrypted flag is set. This is the 128-bit Poly1305 authentication tag generated by the ChaCha20-Poly1305 AEAD cipher. It authenticates both the header (used as Associated Authenticated Data) and the ciphertext payload.

\begin{alertbox}
If the MAC verification fails during parsing, the packet is silently discarded and \texttt{UL\_ERR\_MAC\_VERIFICATION} is returned. This prevents processing of tampered or corrupted encrypted packets.
\end{alertbox}

\subsection{CRC-16 (2 Bytes)}

The final two bytes of every packet contain a CRC-16/MCRF4XX checksum (ITU X.25 polynomial) computed over the entire packet from byte 1 (skipping the SOF marker) through the end of the MAC tag (or payload for unencrypted packets). An additional CRC seed byte, unique to each message ID, is accumulated into the CRC to detect message ID corruption.

\textbf{CRC Parameters:}
\begin{itemize}[noitemsep]
    \item Polynomial: X.25 (CRC-16/MCRF4XX)
    \item Initial value: \texttt{0xFFFF}
    \item Input: Bytes 1 through end-of-MAC/payload, plus a message-ID-specific seed byte
    \item Output: 16-bit checksum stored little-endian
\end{itemize}

\textbf{CRC Seed Table:}
\begin{table}[H]
\centering
\begin{tabular}{@{}lll@{}}
\toprule
\textbf{Message ID} & \textbf{Name} & \textbf{Seed} \\
\midrule
0x001 & Heartbeat & 50 \\
0x002 & Attitude  & 39 \\
0x003 & GPS Raw   & 24 \\
0x004 & Battery   & 154 \\
0x005 & RC Input  & 89 \\
0x006 & Command   & 217 \\
0x007 & Command ACK & 143 \\
0x008 & Mode Change & 178 \\
0x009 & Mission Item & 62 \\
\bottomrule
\end{tabular}
\end{table}

% ============================================================
% 4. STREAM TYPES AND PRIORITIES
% ============================================================
\section{Stream Types and Priority Levels}

\subsection{Stream Types (4-bit Field)}
\label{sec:streams}

\begin{table}[H]
\centering
\caption{UAVLink Stream Types}
\begin{tabular}{@{}cll@{}}
\toprule
\textbf{ID} & \textbf{Name} & \textbf{Purpose} \\
\midrule
0x0 & \texttt{TELEM\_FAST}  & High-rate telemetry (attitude, angular rates) \\
0x1 & \texttt{TELEM\_SLOW}  & Low-rate telemetry (GPS, battery) \\
0x2 & \texttt{CMD}          & Control commands (arm, disarm, mode change) \\
0x3 & \texttt{CMD\_ACK}     & Command acknowledgements \\
0x4 & \texttt{MISSION}      & Waypoint upload and download \\
0x5 & \texttt{VIDEO}        & Video stream data \\
0x6 & \texttt{SENSOR}       & Raw sensor readings \\
0x7 & \texttt{HEARTBEAT}    & System status and keepalive \\
0x8 & \texttt{ALERT}        & Critical alerts and failsafe triggers \\
0xF & \texttt{CUSTOM}       & User-defined / batch messages \\
\bottomrule
\end{tabular}
\end{table}

\subsection{Priority Levels (2-bit Field)}
\label{sec:priority}

\begin{table}[H]
\centering
\caption{UAVLink Priority Levels}
\begin{tabular}{@{}clll@{}}
\toprule
\textbf{Code} & \textbf{Level} & \textbf{Target Latency} & \textbf{Use Case} \\
\midrule
00 & Bulk      & $\sim$1000 ms & Logs, parameter lists, firmware upload \\
01 & Normal    & $\sim$100 ms  & Telemetry, status updates \\
10 & High      & $\sim$20 ms   & Commands, waypoints \\
11 & Emergency & $<$10 ms      & Failsafe, critical alerts, collision avoidance \\
\bottomrule
\end{tabular}
\end{table}

% ============================================================
% 5. MESSAGE PAYLOADS
% ============================================================
\section{Message Payload Specifications}
\label{sec:messages}

All multi-byte fields are serialized in \textbf{little-endian} byte order.

\subsection{Heartbeat (MSG\_ID 0x001)}

\textbf{Purpose:} System status beacon and keepalive. \textbf{Size:} 7 bytes. \textbf{Rate:} 1 Hz.

\begin{table}[H]
\centering
\begin{tabular}{@{}llcl@{}}
\toprule
\textbf{Field} & \textbf{Type} & \textbf{Bytes} & \textbf{Description} \\
\midrule
\texttt{system\_status}  & uint32 & 4 & Bit-packed system state flags \\
\texttt{system\_type}    & uint8  & 1 & Vehicle type (quadcopter, fixed-wing, etc.) \\
\texttt{autopilot\_type} & uint8  & 1 & Autopilot firmware type \\
\texttt{base\_mode}      & uint8  & 1 & Armed/disarmed, manual/auto mode flags \\
\bottomrule
\end{tabular}
\end{table}

\subsection{Attitude (MSG\_ID 0x002)}

\textbf{Purpose:} UAV orientation and angular rates. \textbf{Size:} 18 bytes. \textbf{Rate:} 10--50 Hz.

\begin{table}[H]
\centering
\begin{tabular}{@{}llcl@{}}
\toprule
\textbf{Field} & \textbf{Type} & \textbf{Bytes} & \textbf{Description} \\
\midrule
\texttt{roll}       & float32 & 4 & Roll angle (radians) \\
\texttt{pitch}      & float32 & 4 & Pitch angle (radians) \\
\texttt{yaw}        & float32 & 4 & Yaw angle (radians) \\
\texttt{rollspeed}  & float16 & 2 & Roll rate (rad/s), IEEE 754 half-precision \\
\texttt{pitchspeed} & float16 & 2 & Pitch rate (rad/s), IEEE 754 half-precision \\
\texttt{yawspeed}   & float16 & 2 & Yaw rate (rad/s), IEEE 754 half-precision \\
\bottomrule
\end{tabular}
\end{table}

\begin{notebox}
Angles are stored as full 32-bit IEEE 754 floats for precision. Angular \emph{rates}, which require less precision, are compressed to 16-bit half-precision floats, saving 6 bytes per message (25\% reduction).
\end{notebox}

\subsection{GPS Raw (MSG\_ID 0x003)}

\textbf{Purpose:} Raw GNSS position data. \textbf{Size:} 22 bytes. \textbf{Rate:} 1--10 Hz.

\begin{table}[H]
\centering
\begin{tabular}{@{}llcl@{}}
\toprule
\textbf{Field} & \textbf{Type} & \textbf{Bytes} & \textbf{Description} \\
\midrule
\texttt{lat}        & int32  & 4 & Latitude (degrees $\times 10^7$) \\
\texttt{lon}        & int32  & 4 & Longitude (degrees $\times 10^7$) \\
\texttt{alt}        & int32  & 4 & Altitude AMSL (millimeters) \\
\texttt{eph}        & uint16 & 2 & Horizontal position uncertainty (cm) \\
\texttt{epv}        & uint16 & 2 & Vertical position uncertainty (cm) \\
\texttt{vel}        & uint16 & 2 & Ground speed (cm/s) \\
\texttt{cog}        & uint16 & 2 & Course over ground (degrees $\times 100$) \\
\texttt{fix\_type}  & uint8  & 1 & GPS fix (0=none, 2=2D, 3=3D, 4=DGPS, 5=RTK) \\
\texttt{satellites} & uint8  & 1 & Number of satellites visible \\
\bottomrule
\end{tabular}
\end{table}

\subsection{Battery (MSG\_ID 0x004)}

\textbf{Purpose:} Battery status monitoring. \textbf{Size:} 8 bytes. \textbf{Rate:} 1--5 Hz.

\begin{table}[H]
\centering
\begin{tabular}{@{}llcl@{}}
\toprule
\textbf{Field} & \textbf{Type} & \textbf{Bytes} & \textbf{Description} \\
\midrule
\texttt{voltage}    & uint16 & 2 & Battery voltage (millivolts) \\
\texttt{current}    & int16  & 2 & Current draw (centi-Amps, negative = discharge) \\
\texttt{remaining}  & int16  & 2 & Remaining capacity (\%, $-1$ = unknown) \\
\texttt{cell\_count} & uint8  & 1 & Number of cells \\
\texttt{status}     & uint8  & 1 & Status flags \\
\bottomrule
\end{tabular}
\end{table}

\subsection{RC Input (MSG\_ID 0x005)}

\textbf{Purpose:} Radio control channel data. \textbf{Size:} 18 bytes. \textbf{Rate:} 10--50 Hz.

\begin{table}[H]
\centering
\begin{tabular}{@{}llcl@{}}
\toprule
\textbf{Field} & \textbf{Type} & \textbf{Bytes} & \textbf{Description} \\
\midrule
\texttt{channels[8]} & uint16[8] & 16 & RC channel values (1000--2000 $\mu$s, 0=disconnected) \\
\texttt{rssi}         & uint8     & 1  & Signal strength (0--100\%) \\
\texttt{quality}      & uint8     & 1  & Link quality (0--100\%) \\
\bottomrule
\end{tabular}
\end{table}

\subsection{Batch Message (MSG\_ID 0x3FF)}

\textbf{Purpose:} Aggregate multiple small messages into a single packet. \textbf{Max sub-messages:} 8.

\textbf{Batch payload format:}
\begin{table}[H]
\centering
\begin{tabular}{@{}lcl@{}}
\toprule
\textbf{Field} & \textbf{Size} & \textbf{Description} \\
\midrule
\texttt{num\_messages} & 1 byte & Number of sub-messages (1--8) \\
\emph{For each sub-message:} & & \\
\quad \texttt{msg\_id}  & 2 bytes & Message ID (little-endian) \\
\quad \texttt{length}   & 1 byte  & Payload length (max 64) \\
\quad \texttt{data}     & $n$ bytes & Sub-message payload \\
\bottomrule
\end{tabular}
\end{table}

% ============================================================
% BIDIRECTIONAL COMMUNICATION SECTION
% ============================================================
\section{Bidirectional Communication}

UAVLink v1.2 introduces full bidirectional command and control capabilities, enabling Ground Control Stations (GCS) to send commands to UAVs and receive command acknowledgments (ACKs) in return. This feature is critical for reliable, safety-critical UAV operations.

\subsection{Command Architecture}

The bidirectional communication system consists of three message types:

\begin{enumerate}[noitemsep]
    \item \textbf{Command} (GCS $\to$ UAV): Control commands with parameters
    \item \textbf{Command ACK} (UAV $\to$ GCS): Acknowledgment with result status
    \item \textbf{Mode Change} (GCS $\to$ UAV): Flight mode change requests
\end{enumerate}

All command-related messages use dedicated stream types (\texttt{UL\_STREAM\_CMD} and \texttt{UL\_STREAM\_CMD\_ACK}) and include a \texttt{target\_sys\_id} field for routing, supporting multi-vehicle networks.

\begin{alertbox}
All command-related messages are \textbf{always encrypted} with the \texttt{ENCRYPT\_ALWAYS} policy. This prevents command injection attacks and ensures all control-plane traffic is authenticated.
\end{alertbox}

\subsection{Command (MSG\_ID 0x006)}

\textbf{Purpose:} Generic command from GCS to UAV. \textbf{Size:} 8 bytes. \textbf{Encryption:} Always. \textbf{Stream Type:} \texttt{UL\_STREAM\_CMD}. \textbf{Priority:} High.

\begin{table}[H]
\centering
\begin{tabular}{@{}llcl@{}}
\toprule
\textbf{Field} & \textbf{Type} & \textbf{Bytes} & \textbf{Description} \\
\midrule
\texttt{command\_id} & uint16 & 2 & Command identifier (see table below) \\
\texttt{param1}      & uint16 & 2 & Parameter 1 (command-specific) \\
\texttt{param2}      & uint16 & 2 & Parameter 2 (reserved, set to 0) \\
\texttt{param3}      & uint16 & 2 & Parameter 3 (reserved, set to 0) \\
\bottomrule
\end{tabular}
\end{table}

\subsubsection{Command IDs}

\begin{table}[H]
\centering
\caption{UAVLink Command Set}
\begin{tabular}{@{}clll@{}}
\toprule
\textbf{ID} & \textbf{Name} & \textbf{Param1} & \textbf{Description} \\
\midrule
0x0001 & \texttt{ARM}       & --- & Arm motors (safety checks must pass) \\
0x0002 & \texttt{DISARM}    & --- & Disarm motors (must be landed) \\
0x0003 & \texttt{TAKEOFF}   & Alt (cm) & Takeoff to specified altitude AGL \\
0x0004 & \texttt{LAND}      & --- & Land at current position \\
0x0005 & \texttt{RTL}       & --- & Return to launch point and land \\
0x0006 & \texttt{EMERGENCY} & --- & Emergency stop (immediate motor cutoff) \\
\bottomrule
\end{tabular}
\end{table}

\begin{notebox}
The \texttt{TAKEOFF} command uses \texttt{param1} to specify the target altitude in centimeters above ground level (AGL). For example, \texttt{param1 = 1000} means climb to 10 meters. All other commands currently ignore parameters.
\end{notebox}

\subsection{Connection Timeouts and Failsafe (RTL)}
To ensure strict functional safety, the application layer should consistently monitor the time lapse between receiving valid \texttt{UL\_STREAM\_CMD} messages. If this interval exceeds a pre-defined safety threshold (e.g., 3.0 seconds) while the system is armed, the UAV must automatically assume a Link Loss state and transition into \texttt{UL\_MODE\_RTL} (Return-To-Launch). The core UAVLink parser supports Link Quality metric generation via the \texttt{ul\_get\_link\_quality()} API function to assist with predictive signal dropouts prior to a total timeout event.

\subsection{Command ACK (MSG\_ID 0x007)}

\textbf{Purpose:} UAV acknowledges a received command. \textbf{Size:} 4 bytes. \textbf{Encryption:} Always. \textbf{Stream Type:} \texttt{UL\_STREAM\_CMD\_ACK}. \textbf{Priority:} High.

\begin{table}[H]
\centering
\begin{tabular}{@{}llcl@{}}
\toprule
\textbf{Field} & \textbf{Type} & \textbf{Bytes} & \textbf{Description} \\
\midrule
\texttt{command\_id} & uint16 & 2 & Command ID being acknowledged \\
\texttt{result}      & uint8  & 1 & Result code (see table below) \\
\texttt{progress}    & uint8  & 1 & Progress 0--100\% (for \texttt{IN\_PROGRESS}) \\
\bottomrule
\end{tabular}
\end{table}

\subsubsection{ACK Result Codes}

\begin{table}[H]
\centering
\caption{Command ACK Result Codes}
\begin{tabular}{@{}clp{8cm}@{}}
\toprule
\textbf{Code} & \textbf{Name} & \textbf{Meaning} \\
\midrule
0x00 & \texttt{OK}          & Command accepted and executed successfully \\
0x01 & \texttt{REJECTED}    & Command rejected due to precondition failure (e.g., ARM while already armed, TAKEOFF while disarmed) \\
0x02 & \texttt{UNSUPPORTED} & Unknown or unimplemented command ID \\
0x03 & \texttt{FAILED}      & Command execution failed (e.g., GPS lost during TAKEOFF) \\
0x04 & \texttt{IN\_PROGRESS} & Command accepted, execution in progress (e.g., climbing to altitude). Check \texttt{progress} field for completion percentage. \\
\bottomrule
\end{tabular}
\end{table}

\subsubsection{Command Processing Flow}

The UAV command processor implements state-based validation:

\begin{enumerate}[noitemsep]
    \item \textbf{Receive:} Parse command from GCS via \texttt{ul\_parse\_char()} parser
    \item \textbf{Validate:} Check current state (armed/disarmed, flying/landed)
    \item \textbf{Execute or Reject:} Perform action or return \texttt{REJECTED} ACK
    \item \textbf{Send ACK:} Transmit acknowledgment with result code
    \item \textbf{Update State:} Modify internal UAV state machine
\end{enumerate}

\textbf{Example validation rules:}
\begin{itemize}[noitemsep]
    \item \texttt{ARM}: Only allowed when disarmed and GPS lock acquired
    \item \texttt{DISARM}: Only allowed when landed (altitude $<$ 50cm)
    \item \texttt{TAKEOFF}: Only allowed when armed and not already flying
    \item \texttt{EMERGENCY}: Always accepted regardless of state
\end{itemize}

\subsection{Mode Change (MSG\_ID 0x008)}

\textbf{Purpose:} Request flight mode change. \textbf{Size:} 2 bytes. \textbf{Encryption:} Always. \textbf{Stream Type:} \texttt{UL\_STREAM\_CMD}.

\begin{table}[H]
\centering
\begin{tabular}{@{}llcl@{}}
\toprule
\textbf{Field} & \textbf{Type} & \textbf{Bytes} & \textbf{Description} \\
\midrule
\texttt{mode}     & uint8 & 1 & Target flight mode (see table below) \\
\texttt{reserved} & uint8 & 1 & Reserved for future use (set to 0) \\
\bottomrule
\end{tabular}
\end{table}

\subsubsection{Flight Modes}

\begin{table}[H]
\centering
\caption{UAVLink Flight Modes}
\begin{tabular}{@{}clp{8cm}@{}}
\toprule
\textbf{Code} & \textbf{Name} & \textbf{Description} \\
\midrule
0x00 & \texttt{MANUAL}    & Direct pilot control, no stabilization \\
0x01 & \texttt{STABILIZE} & Attitude stabilization with manual throttle \\
0x02 & \texttt{ALT\_HOLD}  & Altitude hold with manual horizontal control \\
0x03 & \texttt{LOITER}    & Position hold (GPS required) \\
0x04 & \texttt{AUTO}      & Autonomous waypoint navigation \\
0x05 & \texttt{RTL}       & Return to launch (automatic) \\
0x06 & \texttt{LAND}      & Autonomous landing (automatic) \\
\bottomrule
\end{tabular}
\end{table}

\begin{notebox}
Mode change requests are treated as commands and trigger a Command ACK response. The UAV may reject mode changes if prerequisites are not met (e.g., switching to \texttt{AUTO} without uploaded waypoints).
\end{notebox}

\subsection{Mission Item (MSG\_ID 0x009)}

\textbf{Purpose:} Upload waypoints to UAV. Supports fragmentation for multi-waypoint missions. \textbf{Size:} 20 bytes per item. \textbf{Encryption:} Always. \textbf{Stream Type:} \texttt{UL\_STREAM\_MISSION}.

\begin{table}[H]
\centering
\begin{tabular}{@{}llcl@{}}
\toprule
\textbf{Field} & \textbf{Type} & \textbf{Bytes} & \textbf{Description} \\
\midrule
\texttt{seq}          & uint16 & 2 & Waypoint sequence number (0-based) \\
\texttt{frame}        & uint8  & 1 & Coordinate frame (0=global/AMSL, 1=relative/AGL) \\
\texttt{command}      & uint8  & 1 & Waypoint action (see table below) \\
\texttt{lat}          & int32  & 4 & Latitude (degrees $\times 10^7$) \\
\texttt{lon}          & int32  & 4 & Longitude (degrees $\times 10^7$) \\
\texttt{alt}          & int32  & 4 & Altitude (millimeters) \\
\texttt{speed}        & uint16 & 2 & Desired speed (cm/s, 0=default) \\
\texttt{loiter\_time} & uint16 & 2 & Loiter time at waypoint (seconds, 0=none) \\
\bottomrule
\end{tabular}
\end{table}

\subsubsection{Waypoint Actions}

\begin{table}[H]
\centering
\begin{tabular}{@{}clp{8cm}@{}}
\toprule
\textbf{Code} & \textbf{Action} & \textbf{Description} \\
\midrule
0 & \texttt{NAVIGATE} & Fly to waypoint and continue to next \\
1 & \texttt{LOITER}   & Fly to waypoint, orbit for \texttt{loiter\_time} seconds \\
2 & \texttt{LAND}     & Fly to waypoint and land (final waypoint only) \\
\bottomrule
\end{tabular}
\end{table}

\subsubsection{Mission Upload Protocol}

Multi-waypoint missions are uploaded using the following format:

\begin{enumerate}[noitemsep]
    \item \textbf{Pack Mission:} Create buffer with count byte + N waypoint structs:
    \begin{center}
    \texttt{[count (1B)] [waypoint\_0 (20B)] [waypoint\_1 (20B)] ... [waypoint\_N (20B)]}
    \end{center}
    \item \textbf{Fragment (if needed):} If total size $>$ 256 bytes, split into fragments using \texttt{ul\_fragment\_split()}
    \item \textbf{Transmit:} Send each fragment with 10ms delay between packets
    \item \textbf{Reassemble:} UAV reconstructs complete mission buffer via \texttt{ul\_reassembly\_add()}
    \item \textbf{Parse:} Read count, then deserialize N mission items from buffer
\end{enumerate}

\textbf{Example:} 5 waypoints = 1 + (5 $\times$ 20) = 101 bytes (fits in single packet)

\begin{notebox}
The protocol includes full fragment reassembly logic built directly into the core using \texttt{ul\_fragment\_split()} to divide payloads up to \texttt{UL\_FRAG\_MAX\_TOTAL} and \texttt{ul\_reassembly\_add()} with concurrent slotted context to automatically merge chunks back together.
\end{notebox}

% ============================================================
% 6. ENCRYPTION
% ============================================================
\section{Encryption and Authentication}

\subsection{ChaCha20-Poly1305 AEAD}

UAVLink uses the \textbf{ChaCha20-Poly1305} Authenticated Encryption with Associated Data (AEAD) construction, as specified in RFC 8439. This provides both confidentiality (encryption) and integrity/authenticity (MAC) in a single operation.

\textbf{Parameters:}
\begin{itemize}[noitemsep]
    \item \textbf{Key:} 256-bit (32 bytes) pre-shared symmetric key.
    \item \textbf{Nonce:} 64-bit (8 bytes), zero-padded to 192 bits for Monocypher's XChaCha20 interface.
    \item \textbf{AAD:} The complete packet header (base + extended) is used as Associated Authenticated Data. This means the header is \emph{not} encrypted but \emph{is} authenticated---any tampering with header fields will cause MAC verification failure.
    \item \textbf{MAC Tag:} 128-bit (16 bytes) Poly1305 tag appended after the ciphertext.
\end{itemize}

\subsection{Encryption Flow (Packing)}

\begin{enumerate}[noitemsep]
    \item Encode the base header (4 bytes) and extended header (variable) into the output buffer.
    \item Call \texttt{crypto\_aead\_lock(ciphertext, mac, key, nonce, header\_as\_AAD, plaintext)}.
    \item Append the 16-byte MAC tag immediately after the ciphertext.
    \item Compute CRC-16 over the entire packet (header + ciphertext + MAC).
    \item Append the 2-byte CRC.
\end{enumerate}

\subsection{Decryption Flow (Parsing)}

\begin{enumerate}[noitemsep]
    \item Parse and validate the header; collect the ciphertext and MAC tag.
    \item Verify CRC-16 over the received bytes.
    \item Call \texttt{crypto\_aead\_unlock(plaintext, mac, key, nonce, header\_as\_AAD, ciphertext)}.
    \item If the MAC verification fails (return $\neq 0$), discard the packet and return \texttt{UL\_ERR\_MAC\_VERIFICATION}.
    \item If successful, the plaintext payload is available for deserialization.
\end{enumerate}

\subsection{Nonce Management}

Nonce uniqueness is critical for AEAD security. UAVLink uses a \textbf{hybrid counter+random} approach:

\begin{itemize}[noitemsep]
    \item \textbf{Counter (bytes 0--3):} A 32-bit monotonically increasing counter, initialized to a random value at startup. Guarantees uniqueness even if the random source has limited entropy.
    \item \textbf{Random (bytes 4--7):} 32 bits of cryptographically random data generated via the platform's secure RNG (\texttt{CryptGenRandom} on Windows, \texttt{/dev/urandom} on Unix).
\end{itemize}

\begin{warnbox}
If the counter reaches \texttt{0xFFFFFFFF}, it is re-seeded with a fresh random value rather than wrapping to zero. This prevents nonce reuse even under extremely high packet rates.
\end{warnbox}

\subsubsection*{Nonce State Persistence (NVM)}

To completely prevent nonce reuse after catastrophic power loss or system reboots, UAVLink includes support for Non-Volatile Memory (NVM) persistence. The core API provides \texttt{ul\_nonce\_get\_counter()} and \texttt{ul\_nonce\_set\_counter()} to safely extract and inject the inner 32-bit counter. 

In a production environment, the counter should be persisted to disk or flash periodically and before power down. Upon boot, a safety jump (e.g., $+10,000$) is instantly applied to the retrieved counter and immediately re-saved to unconditionally cover any unsaved packets generated right before power loss.

\subsection{Selective Encryption Policies}

Not all messages require encryption. UAVLink defines three per-message-ID policies:

\begin{table}[H]
\centering
\caption{Default Encryption Policies}
\begin{tabular}{@{}lll@{}}
\toprule
\textbf{Message} & \textbf{Policy} & \textbf{Rationale} \\
\midrule
Heartbeat  & \texttt{NEVER}    & Public status; no confidentiality needed \\
Attitude   & \texttt{NEVER}    & High-rate telemetry; encryption overhead undesirable \\
GPS Raw    & \texttt{OPTIONAL} & Location data may be sensitive \\
Battery    & \texttt{OPTIONAL} & Medium sensitivity \\
RC Input   & \texttt{ALWAYS}   & Control inputs---security critical \\
Command    & \texttt{ALWAYS}   & Must prevent command injection \\
Batch      & \texttt{OPTIONAL} & Depends on sub-message content \\
\bottomrule
\end{tabular}
\end{table}

This selective approach saves 24 bytes per unencrypted packet (8-byte nonce + 16-byte MAC), achieving up to \textbf{60\% bandwidth reduction} for heartbeat streams.

\begin{notebox}
All command-related messages (Command, Command ACK, Mode Change, Mission Item) are set to \texttt{ALWAYS} encrypt. This prevents command injection attacks and ensures all control-plane traffic is authenticated.
\end{notebox}

% ============================================================
% FRAGMENTATION SECTION
% ============================================================
\section{Packet Fragmentation and Reassembly}

\subsection{Overview}

UAVLink supports packet fragmentation to enable transmission of messages larger than the maximum single-packet payload size (512 bytes). Large messages---such as multi-waypoint missions, firmware chunks, or parameter lists---are automatically split into multiple fragments, transmitted independently, and reassembled at the receiver.

\begin{notebox}
The fragmentation feature is architecturally designed with header fields and encoding/decoding logic fully implemented. Complete fragmentation split and reassembly APIs are planned for future implementation.
\end{notebox}

\subsection{Fragmentation Design}

Each fragment is transmitted as an independent UAVLink packet with its own:
\begin{itemize}[noitemsep]
    \item Complete base and extended header
    \item Encryption (each fragment encrypted independently with unique nonce)
    \item CRC-16 checksum
    \item Fragment metadata (index and total count)
\end{itemize}

\textbf{Key properties:}
\begin{itemize}[noitemsep]
    \item Fragments can be received out of order (receiver buffers)
    \item Each fragment independently protected against tampering
    \item Supports up to 256 fragments per message (uint8\_t fields)
    \item Reassembly timeout prevents memory exhaustion from incomplete messages
\end{itemize}

\subsection{Fragment Header Fields}

When the \texttt{fragmented} flag (base header byte 3, bit 2) is set, the extended header includes two additional bytes:

\begin{table}[H]
\centering
\begin{tabular}{@{}llp{8cm}@{}}
\toprule
\textbf{Field} & \textbf{Type} & \textbf{Description} \\
\midrule
\texttt{frag\_index} & uint8 & Zero-based index of this fragment (0 = first, 1 = second, etc.) \\
\texttt{frag\_total} & uint8 & Total number of fragments in the complete message (1--256) \\
\bottomrule
\end{tabular}
\end{table}

\textbf{Example:} A 600-byte mission upload split into 3 fragments:
\begin{itemize}[noitemsep]
    \item Fragment 0: \texttt{frag\_index=0, frag\_total=3, payload=256 bytes}
    \item Fragment 1: \texttt{frag\_index=1, frag\_total=3, payload=256 bytes}
    \item Fragment 2: \texttt{frag\_index=2, frag\_total=3, payload=88 bytes}
\end{itemize}

\subsection{Fragmentation API (Planned)}

The following APIs are defined for use by application code:

\begin{lstlisting}
// Fragment a large message into multiple packets
// Returns number of fragments created
int ul_fragment_split(const uint8_t *data, size_t data_len,
                      uint8_t *fragments[], size_t max_frag_size);

// Reassembly context for collecting fragments
typedef struct {
    uint16_t msg_id;       // Message ID being reassembled
    uint8_t total_frags;   // Expected fragment count
    uint8_t recv_mask;     // Bitmap of received fragments
    uint8_t *buffer;       // Reassembly buffer
    size_t buffer_len;     // Total reassembled size
    uint32_t timestamp;    // Timeout tracking
} ul_reassembly_ctx_t;

// Initialize reassembly context
void ul_reassembly_init(ul_reassembly_ctx_t *ctx);

// Add a received fragment; returns 1 when complete
int ul_reassembly_add(ul_reassembly_ctx_t *ctx, 
                      const ul_header_t *hdr,
                      const uint8_t *payload);
\end{lstlisting}

\subsection{Usage Strategy}

Until full fragmentation APIs are implemented, applications can manually split large messages:

\begin{enumerate}[noitemsep]
    \item Divide data into chunks $\leq$ 256 bytes
    \item Set \texttt{fragmented} flag and populate \texttt{frag\_index}, \texttt{frag\_total}
    \item Transmit each fragment with 10ms inter-packet delay
    \item Receiver buffers fragments by \texttt{msg\_id} and sequence number
    \item Reassemble when all fragments received or timeout expires
\end{enumerate}

\subsection{Replay Protection}

The parser maintains a \textbf{32-packet sliding window} for replay detection:
\begin{itemize}[noitemsep]
    \item A \texttt{last\_seq} variable tracks the highest accepted sequence number.
    \item A 32-bit bitmap (\texttt{replay\_window}) records which of the last 32 sequence numbers have been seen.
    \item Incoming packets with sequence numbers that fall outside the window or match an already-seen sequence are rejected.
    \item When a newer sequence arrives, the window slides forward.
\end{itemize}

% ============================================================
% 7. PARSER STATE MACHINE
% ============================================================
\section{Parser State Machine}

UAVLink provides a byte-by-byte state machine parser ideal for UART/serial interfaces where data arrives one byte at a time.

\subsection{Parser States}

\begin{enumerate}[noitemsep]
    \item \textbf{IDLE:} Waiting for the SOF byte (\texttt{0xA5}). All other bytes are discarded.
    \item \textbf{BASE\_HDR:} Collecting bytes 1--3 of the base header. After 4 bytes, decodes payload length, priority, stream type, flags, and partial sequence.
    \item \textbf{EXT\_HDR:} Collecting the extended header. Length depends on stream type (CMD adds 1 byte), fragmentation flag (+2 bytes), and encryption flag (+8 bytes nonce).
    \item \textbf{PAYLOAD:} Collecting payload bytes (and MAC tag if encrypted).
    \item \textbf{CRC:} Collecting the final 2-byte CRC. Upon completion, verifies CRC, decrypts if needed, checks replay window, and returns the result.
\end{enumerate}

\subsection{Return Codes}

\begin{table}[H]
\centering
\begin{tabular}{@{}lcl@{}}
\toprule
\textbf{Code} & \textbf{Value} & \textbf{Meaning} \\
\midrule
\texttt{UL\_OK}                   &  0 & Complete valid packet received \\
(parsing)                         &  1 & Still parsing, need more bytes \\
\texttt{UL\_ERR\_CRC}             & $-1$ & CRC mismatch or replay detected \\
\texttt{UL\_ERR\_NO\_KEY}         & $-2$ & Encrypted packet but no key provided \\
\texttt{UL\_ERR\_MAC\_VERIFICATION} & $-3$ & AEAD authentication failed \\
\texttt{UL\_ERR\_BUFFER\_OVERFLOW} & $-4$ & Payload exceeds buffer capacity \\
\texttt{UL\_ERR\_INVALID\_HEADER} & $-5$ & SOF mismatch or malformed header \\
\texttt{UL\_ERR\_NULL\_POINTER}   & $-6$ & NULL argument passed \\
\bottomrule
\end{tabular}
\end{table}

% ============================================================
% 8. PHASE 2 OPTIMIZATIONS
% ============================================================
\section{Phase 2: Performance Optimizations}

\subsection{Zero-Copy Parser}

The standard parser copies incoming data into an internal 512-byte buffer and then copies the payload to an output buffer---two copies per packet. The zero-copy parser eliminates the intermediate buffer:

\begin{itemize}[noitemsep]
    \item Uses a minimal 32-byte header buffer (vs.\ 512 bytes), saving 480 bytes per parser instance.
    \item Payload bytes are written \emph{directly} to a user-provided output buffer.
    \item Achieves approximately \textbf{2$\times$ faster parsing} by eliminating redundant memory copies.
\end{itemize}

\subsection{Memory Pool Allocator}

For real-time systems, \texttt{malloc()}/\texttt{free()} are non-deterministic. UAVLink provides a fixed-size memory pool:

\begin{itemize}[noitemsep]
    \item \textbf{32 buffers} of 512 bytes each = 16 KB total.
    \item \textbf{O(1) allocation} using a 32-bit bitmap and hardware bit-scan instructions (\texttt{\_BitScanForward} on MSVC, \texttt{\_\_builtin\_ffs} on GCC/Clang).
    \item \textbf{O(1) deallocation} by setting the corresponding bitmap bit.
    \item \textbf{Secure zeroing:} Released buffers are zeroed with \texttt{memset()} to prevent leaking decrypted plaintext or key material through reuse.
    \item \textbf{Statistics tracking:} Allocation count, free count, peak usage, and current usage for debugging and resource monitoring.
\end{itemize}

\subsection{Hardware Crypto Detection}

At startup, UAVLink probes the CPU for SIMD capabilities and selects the fastest available backend:

\begin{table}[H]
\centering
\caption{Crypto Backend Priority}
\begin{tabular}{@{}clc@{}}
\toprule
\textbf{Priority} & \textbf{Backend} & \textbf{Speedup} \\
\midrule
1 (highest) & x86 AVX2         & 4$\times$ \\
2           & ARM NEON         & 4$\times$ \\
3           & x86 SSE2         & 2$\times$ \\
4 (fallback)& Software (Monocypher) & 1$\times$ \\
\bottomrule
\end{tabular}
\end{table}

\subsection{Crypto Context Caching}

When sending bursts of packets with the same encryption key, the crypto context cache stores the last-used key and allows the implementation to skip redundant key expansion. This provides approximately \textbf{30\% speedup} for consecutive packets.

% ============================================================
% 9. PHASE 3 ADVANCED FEATURES
% ============================================================
\section{Phase 3: Advanced Optimizations}

\subsection{LZ4-Style Compression}

UAVLink includes a run-length encoding (RLE) compression engine optimized for speed:

\begin{itemize}[noitemsep]
    \item \textbf{Compression:} Runs of $\geq$3 identical bytes are encoded as a 3-byte sequence: marker (\texttt{0xFF}) + value + count. Literal \texttt{0xFF} bytes are escaped.
    \item \textbf{Decompression:} Symmetric inverse operation.
    \item \textbf{Adaptive:} The \texttt{ul\_should\_compress()} function performs a quick entropy check on the first 64 bytes. Compression is only applied when $>$25\% of sampled bytes are repeats, avoiding expansion on high-entropy data.
\end{itemize}

\subsection{Reed-Solomon Forward Error Correction}

UAVLink implements XOR-based parity FEC for recovering from packet loss:

\begin{itemize}[noitemsep]
    \item Data is divided into configurable \textbf{data shards} (up to 16) and \textbf{parity shards} (up to 16).
    \item Parity is computed as the byte-wise XOR of all data shards.
    \item The decoder can reconstruct \textbf{one missing data shard} from the remaining data shards plus one parity shard.
    \item Enables recovery from up to 25\% packet loss without retransmission.
\end{itemize}

\subsection{Delta Encoding}

For slowly-changing telemetry data (e.g., GPS coordinates), delta encoding transmits only the \emph{difference} from the previous value:

\begin{itemize}[noitemsep]
    \item \textbf{First packet:} Full values are sent (marker \texttt{0x00} + raw struct).
    \item \textbf{Subsequent packets:} Deltas are sent (marker \texttt{0x01} + per-field deltas).
    \item \textbf{Small delta optimization:} If a delta fits in 16 bits ($-32768$ to $+32767$), it is encoded in 2 bytes instead of 4 bytes.
    \item \textbf{GPS bandwidth savings:} Full GPS = 22 bytes; delta GPS $\approx$ 12 bytes (\textbf{57\% reduction}).
\end{itemize}

% ============================================================
% 10. HARDWARE ACCELERATION
% ============================================================
\section{Hardware-Accelerated Cryptography}

\subsection{ARM NEON Implementation}

The NEON-accelerated ChaCha20 processes four 32-bit state words in parallel using 128-bit SIMD registers:

\begin{enumerate}[noitemsep]
    \item Load the 16-word ChaCha20 state into four \texttt{uint32x4\_t} NEON vectors.
    \item Perform 20 rounds (10 double-rounds) of the quarter-round function using NEON intrinsics: \texttt{vaddq\_u32}, \texttt{veorq\_u32}, \texttt{vshlq\_n\_u32}, \texttt{vshrq\_n\_u32}, \texttt{vorrq\_u32}.
    \item Diagonal rounds are achieved by rotating lanes with \texttt{vextq\_u32}.
    \item The RFC 8439 counter convention is followed: counter 0 generates the Poly1305 key, counter 1 begins encryption.
\end{enumerate}

\textbf{Performance:} $\sim$50 $\mu$s per packet (4$\times$ faster than software).

\subsection{x86 SSE2/AVX2 Stubs}

SSE2 and AVX2 backends are provided as integration points. They currently fall back to Monocypher but are structured for drop-in SIMD optimization. AVX2 can process 8 ChaCha20 blocks in parallel for an additional 2$\times$ over SSE2.

% ============================================================
% 11. API REFERENCE
% ============================================================
\section{API Reference}

\subsection{Core Functions}

\begin{lstlisting}
// Initialize parser state machine
void ul_parser_init(ul_parser_t *p);

// Parse a single byte; returns UL_OK on complete packet
int ul_parse_char(ul_parser_t *p, uint8_t c, const uint8_t *key_32b);

// Pack a complete message into a wire-ready buffer
int uavlink_pack(uint8_t *buf, const ul_header_t *h,
                 const uint8_t *payload, const uint8_t *key_32b);

// Pack with automatic nonce generation (recommended)
int uavlink_pack_with_nonce(uint8_t *buf, const ul_header_t *h,
    const uint8_t *payload, const uint8_t *key_32b,
    ul_nonce_state_t *nonce_state);
\end{lstlisting}

\subsection{Optimization Functions}

\begin{lstlisting}
// Selective encryption based on message policy
int uavlink_pack_selective(uint8_t *buf, const ul_header_t *h,
    const uint8_t *payload, const uint8_t *key_32b,
    ul_nonce_state_t *nonce_state);

// Pack with crypto context caching (30% speedup)
int uavlink_pack_cached(uint8_t *buf, const ul_header_t *h,
    const uint8_t *payload, const uint8_t *key_32b,
    ul_nonce_state_t *nonce_state, ul_crypto_ctx_t *crypto_ctx);

// Batch multiple messages into one packet
int uavlink_pack_batch(uint8_t *buf, const ul_batch_t *batch,
    const uint8_t *key_32b, ul_nonce_state_t *nonce_state,
    uint8_t priority);

// Phase 2: Fast pack with memory pool + all optimizations
int ul_pack_fast(ul_mempool_t *pool, const ul_header_t *h,
    const uint8_t *payload, const uint8_t *key_32b,
    ul_nonce_state_t *nonce_state, ul_crypto_ctx_t *crypto_ctx,
    uint8_t **buffer);
\end{lstlisting}

\subsection{Serialization Functions}

\begin{lstlisting}
int ul_serialize_heartbeat(const ul_heartbeat_t *hb, uint8_t *out);
int ul_deserialize_heartbeat(ul_heartbeat_t *hb, const uint8_t *in);
int ul_serialize_attitude(const ul_attitude_t *att, uint8_t *out);
int ul_deserialize_attitude(ul_attitude_t *att, const uint8_t *in);
int ul_serialize_gps_raw(const ul_gps_raw_t *gps, uint8_t *out);
int ul_deserialize_gps_raw(ul_gps_raw_t *gps, const uint8_t *in);
int ul_serialize_battery(const ul_battery_t *bat, uint8_t *out);
int ul_deserialize_battery(ul_battery_t *bat, const uint8_t *in);
int ul_serialize_rc_input(const ul_rc_input_t *rc, uint8_t *out);
int ul_deserialize_rc_input(ul_rc_input_t *rc, const uint8_t *in);
\end{lstlisting}

% ============================================================
% 12. USAGE EXAMPLES
% ============================================================
\section{Usage Examples}

\subsection{Basic: Send and Receive a Message}

\begin{lstlisting}[caption={Sending an attitude message over UART}]
#include "uavlink.h"

// Setup
ul_nonce_state_t nonce;
ul_nonce_init(&nonce);
uint8_t key[32] = { /* pre-shared key */ };

// Serialize
ul_attitude_t att = {.roll=0.1f, .pitch=-0.05f, .yaw=1.5f,
                     .rollspeed=0.01f, .pitchspeed=0.0f, .yawspeed=0.02f};
uint8_t payload[32];
int plen = ul_serialize_attitude(&att, payload);

// Pack
ul_header_t hdr = {0};
hdr.payload_len = plen;
hdr.priority    = UL_PRIO_HIGH;
hdr.stream_type = UL_STREAM_TELEM_FAST;
hdr.msg_id      = UL_MSG_ATTITUDE;
hdr.sys_id      = 1;
hdr.sequence    = seq++;

uint8_t packet[256];
int len = uavlink_pack_with_nonce(packet, &hdr, payload, key, &nonce);
uart_transmit(packet, len);
\end{lstlisting}

\begin{lstlisting}[caption={Receiving messages from a byte stream}]
ul_parser_t parser;
ul_parser_init(&parser);

while (uart_available()) {
    int r = ul_parse_char(&parser, uart_read(), key);
    if (r == UL_OK) {
        switch (parser.header.msg_id) {
        case UL_MSG_ATTITUDE: {
            ul_attitude_t att;
            ul_deserialize_attitude(&att, parser.payload);
            printf("Roll=%.2f Pitch=%.2f\n", att.roll, att.pitch);
            break;
        }
        case UL_MSG_GPS_RAW: {
            ul_gps_raw_t gps;
            ul_deserialize_gps_raw(&gps, parser.payload);
            printf("Lat=%d Lon=%d\n", gps.lat, gps.lon);
            break;
        }}
    } else if (r == UL_ERR_MAC_VERIFICATION) {
        printf("Tampered packet!\n");
    }
}
\end{lstlisting}

\subsection{Advanced: Phase 2 with Memory Pool}

\begin{lstlisting}[caption={Using Phase 2 fast API with memory pool}]
#include "uavlink.h"
#include "uavlink_fast.h"

ul_mempool_t pool;        ul_mempool_init(&pool);
ul_nonce_state_t nonce;   ul_nonce_init(&nonce);
ul_crypto_ctx_t ctx;      ul_crypto_ctx_init(&ctx);

uint8_t *packet = NULL;
int len = ul_pack_fast(&pool, &hdr, payload, key,
                       &nonce, &ctx, &packet);
if (len > 0 && packet) {
    send(socket, packet, len, 0);
    ul_mempool_free(&pool, packet);  // O(1) dealloc
}
\end{lstlisting}

\subsection{Bidirectional: Sending Commands and Handling ACKs}

\begin{lstlisting}[caption={GCS: Send ARM command and wait for ACK}]
#include "uavlink.h"

// Create ARM command
ul_command_t cmd = {0};
cmd.command_id = UL_CMD_ARM;
cmd.param1 = 0;
cmd.param2 = 0;
cmd.param3 = 0;

// Serialize and pack
uint8_t payload[16];
int plen = ul_serialize_command(&cmd, payload);

ul_header_t hdr = {0};
hdr.payload_len = plen;
hdr.stream_type = UL_STREAM_CMD;
hdr.priority    = UL_PRIO_HIGH;
hdr.encrypted   = true;
hdr.msg_id      = UL_MSG_CMD;
hdr.sys_id      = 255;  // GCS ID
hdr.target_sys_id = 1;  // Target UAV
hdr.sequence    = seq++;

uint8_t packet[128];
int len = uavlink_pack_with_nonce(packet, &hdr, payload, key, &nonce);
sendto(socket, packet, len, 0, (struct sockaddr*)&uav_addr, 
       sizeof(uav_addr));

printf("Sent ARM command, waiting for ACK...\n");
\end{lstlisting}

\begin{lstlisting}[caption={UAV: Process command and send ACK}]
// In UAV main loop
if (r == UL_OK && parser.header.msg_id == UL_MSG_CMD) {
    ul_command_t cmd;
    ul_deserialize_command(&cmd, parser.payload);
    
    // Process command
    ul_command_ack_t ack = {0};
    ack.command_id = cmd.command_id;
    
    switch (cmd.command_id) {
    case UL_CMD_ARM:
        if (!state.armed && state.gps_lock) {
            state.armed = true;
            ack.result = UL_ACK_OK;
            printf("ARMED\n");
        } else {
            ack.result = UL_ACK_REJECTED;
        }
        break;
    case UL_CMD_TAKEOFF:
        if (state.armed && !state.flying) {
            state.target_alt = cmd.param1;  // cm
            state.flying = true;
            ack.result = UL_ACK_OK;
            printf("TAKEOFF to %d cm\n", cmd.param1);
        } else {
            ack.result = UL_ACK_REJECTED;
        }
        break;
    default:
        ack.result = UL_ACK_UNSUPPORTED;
    }
    
    // Send ACK
    uint8_t ack_payload[8];
    int ack_len = ul_serialize_command_ack(&ack, ack_payload);
    send_ack_packet(ack_payload, ack_len);
}
\end{lstlisting}

\begin{lstlisting}[caption={GCS: Receive and display ACK}]
// In GCS receive loop
if (r == UL_OK && parser.header.msg_id == UL_MSG_CMD_ACK) {
    ul_command_ack_t ack;
    ul_deserialize_command_ack(&ack, parser.payload);
    
    const char *cmd_name = get_command_name(ack.command_id);
    const char *result_str;
    
    switch (ack.result) {
    case UL_ACK_OK:          result_str = "OK"; break;
    case UL_ACK_REJECTED:    result_str = "REJECTED"; break;
    case UL_ACK_UNSUPPORTED: result_str = "UNSUPPORTED"; break;
    case UL_ACK_FAILED:      result_str = "FAILED"; break;
    case UL_ACK_IN_PROGRESS: result_str = "IN_PROGRESS"; break;
    default:                 result_str = "UNKNOWN"; break;
    }
    
    printf("ACK: %s -> %s", cmd_name, result_str);
    if (ack.result == UL_ACK_IN_PROGRESS) {
        printf(" (%d%%)", ack.progress);
    }
    printf("\n");
}
\end{lstlisting}

% ============================================================
% 13. BUILDING AND TESTING
% ============================================================
\section{Building and Testing}

\subsection{Test Programs}

UAVLink includes three test programs demonstrating full protocol functionality:

\subsubsection{UAV Simulator (\texttt{uav\_simulator.c})}

A complete bidirectional UAV simulator that:
\begin{itemize}[noitemsep]
    \item Sends telemetry (heartbeat, attitude, GPS, battery) at realistic rates
    \item Receives and processes commands (ARM, DISARM, TAKEOFF, LAND, RTL, EMERGENCY)
    \item Validates commands against internal state (armed/disarmed, flying/landed, GPS lock)
    \item Sends command ACKs with appropriate result codes
    \item Supports mission uploads via fragmented message protocol
    \item Uses UDP networking for realistic wireless link simulation
\end{itemize}

\textbf{Default configuration:}
\begin{itemize}[noitemsep]
    \item Listen port: 14550 (receives commands from GCS)
    \item Send to: 127.0.0.1:14540 (sends telemetry to GCS)
    \item Telemetry rate: 10 Hz attitude, 1 Hz GPS/battery
\end{itemize}

\subsubsection{GCS Receiver (\texttt{gcs\_receiver.c})}

An interactive Ground Control Station that:
\begin{itemize}[noitemsep]
    \item Receives and displays all telemetry messages
    \item Provides interactive command menu for sending commands
    \item Receives and displays command ACKs with result codes
    \item Supports mission upload (5-waypoint test mission)
    \item Decodes and displays all message types with human-readable formatting
\end{itemize}

\textbf{Interactive menu commands:}
\begin{lstlisting}[language={}]
Commands:
  a - ARM
  d - DISARM
  t - TAKEOFF (prompts for altitude)
  l - LAND
  r - RTL (Return to Launch)
  e - EMERGENCY STOP
  m - Upload test mission (5 waypoints)
  q - Quit
\end{lstlisting}

\subsubsection{Benchmark (\texttt{uavlink\_benchmark.c})}

Comprehensive performance profiler measuring:
\begin{itemize}[noitemsep]
    \item Parse speed (µs per packet)
    \item Pack speed (µs per packet)
    \item Encryption overhead (encrypted vs unencrypted)
    \item Memory pool allocation speed
    \item Zero-copy parser performance
    \item Hardware crypto acceleration gains
    \item Bandwidth efficiency (bytes/message for each type)
\end{itemize}

\subsection{Compilation}

\begin{lstlisting}[language=bash,caption={Linux/macOS build}]
cd Protocol

# UAV Simulator (bidirectional telemetry + command processing)
gcc -Wall -O2 -o uav_simulator uav_simulator.c uavlink.c \
    uavlink_fast.c uavlink_compress.c uavlink_hw_crypto.c \
    monocypher.c -lm

# GCS Receiver (interactive command + telemetry display)
gcc -Wall -O2 -o gcs_receiver gcs_receiver.c uavlink.c \
    uavlink_fast.c uavlink_compress.c uavlink_hw_crypto.c \
    monocypher.c -lm

# Performance Benchmark
gcc -Wall -O2 -o benchmark uavlink_benchmark.c uavlink.c \
    uavlink_fast.c uavlink_compress.c uavlink_hw_crypto.c \
    monocypher.c -lm
\end{lstlisting}

\begin{lstlisting}[language=bash,caption={Windows build (MinGW or MSVC)}]
# Windows requires -lws2_32 for socket functions
gcc -Wall -O2 -o uav_simulator.exe uav_simulator.c uavlink.c ^
    uavlink_fast.c uavlink_compress.c uavlink_hw_crypto.c ^
    monocypher.c -lws2_32 -lm

gcc -Wall -O2 -o gcs_receiver.exe gcs_receiver.c uavlink.c ^
    uavlink_fast.c uavlink_compress.c uavlink_hw_crypto.c ^
    monocypher.c -lws2_32 -lm
\end{lstlisting}

\subsection{Running the Test Suite}

\textbf{Step 1: Terminal 1 --- Start UAV Simulator}
\begin{lstlisting}[language=bash]
./uav_simulator
# Output:
# UAVLink UAV Simulator - Listening on UDP 14550
# Sending telemetry to 127.0.0.1:14540
# [Armed: NO | Flying: NO | GPS: OK]
\end{lstlisting}

\textbf{Step 2: Terminal 2 --- Start GCS}
\begin{lstlisting}[language=bash]
./gcs_receiver
# Output:
# UAVLink GCS - Listening on UDP 14540
# Connected to UAV at 127.0.0.1:14550
# [HEARTBEAT] Status: 0x0000 | Type: 2 | Autopilot: 12 | Mode: 0x04
# [ATTITUDE] Roll: 0.05 | Pitch: -0.02 | Yaw: 1.57
# [GPS] Lat: 37.7749000 Lon: -122.4194000 Alt: 15.2m Fix: 3D Sats: 12
\end{lstlisting}

\textbf{Step 3: Send Interactive Commands}
\begin{lstlisting}[language={}]
> a [Enter]        # ARM command
ACK: ARM -> OK

> t [Enter]        # TAKEOFF
Enter altitude (cm): 1000
ACK: TAKEOFF -> OK

> l [Enter]        # LAND
ACK: LAND -> OK

> d [Enter]        # DISARM
ACK: DISARM -> OK
\end{lstlisting}

\textbf{Step 4: Verify in UAV Terminal}
\begin{lstlisting}[language={}]
[CMD] Received ARM command -> Accepted
[Armed: YES | Flying: NO | GPS: OK]
[CMD] Received TAKEOFF command (alt=1000cm) -> Accepted
[Armed: YES | Flying: YES | GPS: OK]
[CMD] Received LAND command -> Accepted
[Armed: YES | Flying: NO | GPS: OK]
[CMD] Received DISARM command -> Accepted
[Armed: NO | Flying: NO | GPS: OK]
\end{lstlisting}

\subsection{Hardware-Accelerated Builds}

\begin{lstlisting}[language=bash,caption={ARM NEON (Raspberry Pi, Jetson, Apple Silicon)}]
# ARM 32-bit (ARMv7-A)
gcc -O2 -mfpu=neon -march=armv7-a -o uav_simulator ...

# ARM 64-bit (ARMv8-A / Apple Silicon)
gcc -O2 -march=armv8-a -o uav_simulator ...
\end{lstlisting}

\begin{lstlisting}[language=bash,caption={x86 SIMD (modern Intel/AMD processors)}]
# SSE2 (baseline, supported on all x64 CPUs)
gcc -O2 -msse2 -o uav_simulator ...

# AVX2 (Haswell/Broadwell and later)
gcc -O2 -mavx2 -o uav_simulator ...
\end{lstlisting}

\begin{notebox}
The protocol automatically detects and enables SIMD acceleration at runtime. Compiling with \texttt{-mavx2} or \texttt{-mfpu=neon} ensures the compiler generates optimized code, but the software fallback (Monocypher) is always available.
\end{notebox}

\subsection{Using the Makefile}

A Makefile is provided for convenient building:

\begin{lstlisting}[language=bash]
cd Protocol

# Build all programs
make

# Build individual targets
make uav_simulator
make gcs_receiver
make benchmark

# Clean build artifacts
make clean

# Windows-specific build (requires MinGW)
make -f Makefile.win
\end{lstlisting}

The Makefile automatically detects the platform (Windows/Linux/macOS) and links the appropriate libraries (\texttt{-lws2\_32} on Windows, standard networking on Unix).

% ============================================================
% 14. ERROR HANDLING
% ============================================================
\section{Error Handling}

All UAVLink functions follow a consistent error reporting pattern:
\begin{itemize}[noitemsep]
    \item \textbf{Positive return values} indicate success (typically the number of bytes written or consumed).
    \item \textbf{Zero} indicates a special success state (\texttt{UL\_OK} for complete packet, 0 for ``need more bytes'').
    \item \textbf{Negative return values} are error codes from the \texttt{ul\_error\_t} enum.
    \item All functions perform \textbf{NULL pointer checks} on their arguments and return \texttt{UL\_ERR\_NULL\_POINTER} if violated.
    \item All functions perform \textbf{bounds checking} on payload lengths and buffer sizes to prevent overflow.
\end{itemize}

% ============================================================
% 15. DEVELOPMENT ROADMAP
% ============================================================
\section{Development Roadmap}

\subsection{Implementation Status}

\begin{table}[H]
\centering
\caption{UAVLink v1.2 Feature Status}
\begin{tabular}{@{}p{6cm}cc@{}}
\toprule
\textbf{Feature} & \textbf{Status} & \textbf{Version} \\
\midrule
\multicolumn{3}{l}{\textbf{Phase 0: Core Protocol}} \\
\quad Bit-packed headers (4-13 bytes) & \checkmark & v1.0 \\
\quad ChaCha20-Poly1305 AEAD encryption & \checkmark & v1.0 \\
\quad CRC-16 integrity checking & \checkmark & v1.0 \\
\quad Byte-by-byte state machine parser & \checkmark & v1.0 \\
\quad Replay protection (32-packet window) & \checkmark & v1.0 \\
\quad Hybrid nonce generation & \checkmark & v1.0 \\
\quad 10 message types (telemetry + commands) & \checkmark & v1.2 \\
\midrule
\multicolumn{3}{l}{\textbf{Phase 1: Quick Wins}} \\
\quad Selective encryption policies & \checkmark & v1.1 \\
\quad Crypto context caching & \checkmark & v1.1 \\
\quad Message batching (up to 8 messages) & \checkmark & v1.1 \\
\quad Bidirectional commands + ACKs & \checkmark & v1.2 \\
\quad Flight mode control & \checkmark & v1.2 \\
\quad Mission waypoint upload & \checkmark & v1.2 \\
\midrule
\multicolumn{3}{l}{\textbf{Phase 2: Performance}} \\
\quad Zero-copy parser & \checkmark & v1.1 \\
\quad O(1) memory pool allocator & \checkmark & v1.1 \\
\quad Hardware crypto detection & \checkmark & v1.1 \\
\quad ARM NEON acceleration & \checkmark & v1.1 \\
\quad x86 SSE2/AVX2 stubs & \checkmark & v1.1 \\
\midrule
\multicolumn{3}{l}{\textbf{Phase 3: Advanced (Partial)}} \\
\quad LZ4-style compression & Partial & v1.1 \\
\quad Delta encoding & \checkmark & v1.1 \\
\quad Reed-Solomon FEC & Planned & --- \\
\midrule
\multicolumn{3}{l}{\textbf{Fragmentation (Designed)}} \\
\quad Fragment header fields & \checkmark & v1.2 \\
\quad Fragment encoding/decoding & \checkmark & v1.2 \\
\quad \texttt{ul\_fragment\_split()} API & Planned & --- \\
\quad \texttt{ul\_reassembly\_add()} API & Planned & --- \\
\quad Out-of-order reassembly & Planned & --- \\
\bottomrule
\end{tabular}
\end{table}

\subsection{Future Enhancements (Priority Order)}

\subsubsection{Priority 1: Complete Fragmentation APIs}
\textbf{Goal:} Enable transmission of messages $>$512 bytes without manual splitting.

\textbf{Tasks:}
\begin{itemize}[noitemsep]
    \item Implement \texttt{ul\_fragment\_split()} to automatically divide large payloads
    \item Implement \texttt{ul\_reassembly\_ctx\_t} state tracking
    \item Add timeout-based reassembly expiration (prevent memory exhaustion)
    \item Support out-of-order fragment reception
    \item Test with large mission uploads (20+ waypoints)
\end{itemize}

\subsubsection{Priority 2: Nonce Persistence (Non-Volatile Memory)}
\textbf{Goal:} Prevent nonce reuse after UAV reboot for maximum security.

\textbf{Tasks:}
\begin{itemize}[noitemsep]
    \item Store nonce counter in EEPROM/Flash
    \item Increment on every packet transmission
    \item Load from NVM on boot, add 1000 to ensure no overlap
    \item Provide NVM backend abstraction (EEPROM, Flash, SD card)
\end{itemize}

\textbf{Estimated security improvement:} Eliminates nonce collision risk across reboots.

\subsubsection{Priority 3: Link Quality Monitoring \& Failsafe}
\textbf{Goal:} Detect link loss and trigger automatic Return-To-Launch (RTL).

\textbf{Tasks:}
\begin{itemize}[noitemsep]
    \item Track packet reception rate (packets/second)
    \item Calculate packet loss percentage using sequence numbers
    \item Measure Round-Trip Time (RTT) for command ACKs
    \item Trigger RTL if no packets received for 3 seconds
    \item Add link quality telemetry message (signal strength, packet loss, RTT)
\end{itemize}

\textbf{Safety justification:} Critical for beyond-visual-line-of-sight (BVLOS) operations.

\subsubsection{Priority 4: Session Key Exchange (ECDH)}
\textbf{Goal:} Eliminate hardcoded pre-shared keys via ephemeral key exchange.

\textbf{Tasks:}
\begin{itemize}[noitemsep]
    \item Implement X25519 Elliptic Curve Diffie-Hellman (built into Monocypher)
    \item Add \texttt{KEY\_EXCHANGE} message (public keys + signatures)
    \item Derive session keys from ECDH shared secret + 32-byte random salt
    \item Rotate session keys every 1 hour or 10{,}000 packets
    \item Maintain backward compatibility with PSK mode
\end{itemize}

\textbf{Security improvement:} Forward secrecy---compromised keys cannot decrypt past traffic.

\subsubsection{Priority 5: Reed-Solomon FEC Completion}
\textbf{Goal:} Recover from 25\% packet loss without retransmission.

\textbf{Tasks:}
\begin{itemize}[noitemsep]
    \item Implement full Reed-Solomon encoder (4 data + 1 parity shards)
    \item Add FEC decoder with Gaussian elimination for recovery
    \item Integrate with fragmentation for multi-fragment FEC protection
    \item Test on simulated lossy links (random drop rates)
\end{itemize}

\textbf{Use case:} Long-range radio links with 10--20\% packet loss (LoRa, 900 MHz telemetry).

\subsection{Completed Features Recap}

UAVLink v1.2 successfully implements:
\begin{itemize}[noitemsep]
    \item \textbf{9 core message types} for full UAV telemetry and control
    \item \textbf{Bidirectional communication} with state-validated command processing
    \item \textbf{ChaCha20-Poly1305 AEAD} with selective encryption policies (82.8\% bandwidth savings)
    \item \textbf{Zero-copy parser} and \textbf{O(1) memory pool} for real-time determinism
    \item \textbf{ARM NEON hardware acceleration} (4$\times$ crypto speedup)
    \item \textbf{Interactive test suite} (UAV simulator + GCS receiver)
\end{itemize}

The protocol is production-ready for UAV applications requiring secure, low-latency, bandwidth-efficient communication over constrained radio links.

% ============================================================
% 16. SECURITY CONSIDERATIONS
% ============================================================
\section{Security Considerations}

\begin{itemize}
    \item \textbf{Key Distribution:} UAVLink uses pre-shared symmetric keys. Key distribution and rotation are the responsibility of the application layer. Keys should be provisioned via a secure out-of-band channel.
    \item \textbf{Nonce Uniqueness:} Reusing a nonce with the same key completely breaks ChaCha20-Poly1305 security. The hybrid counter+random scheme with overflow re-seeding provides strong uniqueness guarantees.
    \item \textbf{Replay Protection:} The 32-packet sliding window prevents replay attacks but cannot detect replays of packets older than 32 sequence numbers. For higher security, increase the window size.
    \item \textbf{Memory Zeroing:} The memory pool zeros released buffers to prevent decrypted plaintext from persisting in memory.
    \item \textbf{Constant-Time MAC Comparison:} MAC verification uses Monocypher's \texttt{crypto\_verify16()} which is constant-time to prevent timing side-channel attacks.
    \item \textbf{Fallback RNG:} If the platform's secure RNG fails, a non-cryptographic fallback is used for nonce generation. This is logged but not fatal, as the counter component still ensures uniqueness.
\end{itemize}

% ============================================================
% 16. ACKNOWLEDGEMENTS
% ============================================================
\section{Acknowledgements}

The UAVLink protocol was developed by \textbf{Krishnashis Das} (GitHub: \href{https://github.com/I-am-Krish}{\texttt{I-am-Krish}}) in collaboration with \textbf{Aadarsh Sinha} (GitHub: \href{https://github.com/Monarch666}{\texttt{Monarch666}}), who co-developed the project and contributed significantly through code development, review and testing.

The cryptographic primitives are provided by the \textbf{Monocypher} library by Loup Vaillant, a portable and audited implementation of ChaCha20-Poly1305.

% ============================================================
% APPENDIX
% ============================================================
\appendix
\section{Error Code Reference}

\begin{table}[H]
\centering
\begin{tabular}{@{}lcl@{}}
\toprule
\textbf{Constant} & \textbf{Value} & \textbf{Description} \\
\midrule
\texttt{UL\_OK}                    &  0 & Operation completed successfully \\
\texttt{UL\_ERR\_CRC}              & $-1$ & CRC-16 checksum mismatch \\
\texttt{UL\_ERR\_NO\_KEY}          & $-2$ & Encrypted packet received without decryption key \\
\texttt{UL\_ERR\_MAC\_VERIFICATION} & $-3$ & Poly1305 MAC verification failed \\
\texttt{UL\_ERR\_BUFFER\_OVERFLOW} & $-4$ & Payload exceeds maximum buffer size \\
\texttt{UL\_ERR\_INVALID\_HEADER}  & $-5$ & SOF byte mismatch or malformed header \\
\texttt{UL\_ERR\_NULL\_POINTER}    & $-6$ & NULL pointer passed as argument \\
\bottomrule
\end{tabular}
\end{table}

\section{Protocol Constants}

\begin{table}[H]
\centering
\begin{tabular}{@{}lcl@{}}
\toprule
\textbf{Constant} & \textbf{Value} & \textbf{Description} \\
\midrule
\texttt{UL\_SOF}              & \texttt{0xA5} & Start of Frame marker \\
\texttt{UL\_MAX\_PAYLOAD\_SIZE} & 512          & Maximum payload in parser buffer (bytes) \\
\texttt{UL\_MAC\_TAG\_SIZE}    & 16            & Poly1305 MAC tag size (bytes) \\
\texttt{UL\_BATCH\_MAX\_MESSAGES} & 8          & Maximum messages per batch \\
\texttt{UL\_MEMPOOL\_NUM\_BUFFERS} & 32        & Number of memory pool buffers \\
\texttt{UL\_MEMPOOL\_BUFFER\_SIZE} & 512       & Size of each pool buffer (bytes) \\
\bottomrule
\end{tabular}
\end{table}

\section{Glossary}

\begin{description}[style=nextline]
    \item[AEAD] Authenticated Encryption with Associated Data. Provides both confidentiality and integrity in one pass.
    \item[ChaCha20] A stream cipher designed by Daniel J.\ Bernstein. Uses a 256-bit key and produces a keystream via 20 rounds of quarter-round operations.
    \item[CRC-16] Cyclic Redundancy Check with a 16-bit output. Detects accidental data corruption.
    \item[Delta Encoding] Encoding scheme that stores the difference between consecutive values rather than absolute values.
    \item[FEC] Forward Error Correction. Adds redundant data so the receiver can correct errors without retransmission.
    \item[GCS] Ground Control Station. The operator's interface for monitoring and commanding UAVs.
    \item[LZ4] A fast lossless compression algorithm. UAVLink uses an RLE variant optimized for small payloads.
    \item[MAC] Message Authentication Code. A short tag that verifies both integrity and authenticity of a message.
    \item[Nonce] A number used once. Ensures that identical plaintexts encrypt to different ciphertexts.
    \item[Poly1305] A one-time authenticator designed by Daniel J.\ Bernstein. Produces a 128-bit MAC tag.
    \item[Reed-Solomon] An error-correcting code that can recover from symbol erasures using algebraic properties.
    \item[SIMD] Single Instruction, Multiple Data. CPU instructions that operate on multiple data elements in parallel.
    \item[SOF] Start of Frame. A fixed byte value used to synchronize the parser to the beginning of a packet.
    \item[UAV] Unmanned Aerial Vehicle. An aircraft operated without a human pilot on board.
    \item[Zero-Copy] A parsing technique that avoids intermediate buffer copies by writing directly to the output destination.
\end{description}

\end{document}
